\subsection{Strict morphisms and properness}

Lemma 1. --- Let E and F be topological vector spaces, let u be a continuous linear map from E into F, and let U be a neighborhood of 0 in E. Suppose that u induces a homeomorphism from U onto a closed subset of F. Then the image I of u is closed and u induces a homeomorphism from E onto I.

The set \(u(\mathrm{U})\) is a neighborhood of 0 in I; it is closed in F, hence the subgroup I of F is locally closed at 0, and consequently it is closed (TG, III, p. 7, prop. 4).

Since \(\operatorname{Ker}(u)\) meets U only at 0, one has \(\operatorname{Ker}(u) = \{0\}\) and the map \(u\) is injective. Let V be a closed neighborhood of 0 in E contained in U. Since \(u(\mathring{\mathrm{V}})\) is a neighborhood of 0 in \(u(\mathrm{V})\) , there exists a balanced neighborhood W of 0 in F such that \(u(\mathrm{V}) \cap \mathrm{W} \subset u(\mathring{\mathrm{V}})\) . The set \(u^{-1}(\mathrm{W})\) is balanced, hence connected. It contains 0 and is contained in \(\mathring{\mathrm{V}} \cup (\mathrm{E} - \mathrm{V})\) . Since \(\mathring{\mathrm{V}}\) and E-V are disjoint open subsets of E, the set \(u^{-1}(\mathrm{W})\) is contained in \(\mathring{\mathrm{V}}\) , whence \(\mathrm{W} \cap \mathrm{I} \subset u(\mathring{\mathrm{V}})\) . Consequently \(u(\mathring{\mathrm{V}})\) is a neighborhood of 0 in I. This implies that \(u\) induces a homeomorphism from E onto I.

PROPOSITION 1. --- Let E be a separated locally convex space, F a locally convex space and u a continuous linear map from E into F. The following conditions are equivalent:

\begin{enumerate}
\def\labelenumi{(\roman{enumi})}
\item
  The map u is a strict morphism, its kernel has finite dimension and its image is closed in F;
\item
  There exists a closed neighborhood V of 0 in E such that the restriction of u to V is proper (TG, I, p. 72).
\item
  \(\Longrightarrow\) (ii): Suppose that u satisfies condition (i). Since the kernel of u has finite dimension, it has a topological supplement \(E_{1}\) in E (III, p. 55, prop. 1), and there exists a compact neighborhood C of 0 in \(\operatorname{Ker}(u)\) . Identify E with \(E_{1} \times \operatorname{Ker}(u)\) ; the set \(V = E_{1} \times C\) is then a closed neighborhood of 0 in E. The restriction of u to V is the composition of the projection of \(E_{1} \times C\) onto \(E_{1}\) which is proper (TG, I, p. 77, cor. 5) and of the restriction \(u_{1}\) of u to \(E_{1}\) . Now \(u_{1}\) is a homeomorphism from \(E_{1}\) onto a closed subspace of F, hence is proper (TG, I, p. 72, prop. 2). The composition of two proper maps is proper (TG, I, p. 73, prop. 5, a)), so the restriction of u to V is proper.
\item
  \(\Longrightarrow\) (i): Let V be a closed neighborhood of 0 in E such that the restriction v of u to V is proper. The set \(\mathrm{V} \cap \mathrm{Ker}(u) = \overline{v}(\{0\})\) is then compact (TG, I, p. 75, th. 1); consequently, the vector space \(\mathrm{Ker}(u)\) is locally compact, hence finite-dimensional (EVT, I, p. 15, th. 3). Let \(E_{1}\) be a topological complement of \(\mathrm{Ker}(u)\) in E (prop. 1 of III, p. 55); set \(V_{1} = E_{1} \cap V\) . The set \(V_{1}\) is closed in V. The map \(u|V_{1}\) is proper (TG, I, p. 74, cor. 1) and injective, hence is a homeomorphism from \(V_{1}\) onto a closed subset of F (TG, I, p. 72, prop. 2). By Lemma 1, the restriction of u to \(E_{1}\) is a homeomorphism from \(E_{1}\) onto a closed subspace of F, hence u is a strict morphism with closed image.
\end{enumerate}

